% Generated draft: requires mathematical and source review.
\documentclass[11pt]{article}
\usepackage[margin=25mm]{geometry}
\usepackage{fontspec}
\setmainfont{Latin Modern Roman}
\usepackage{amsmath,amssymb,mathtools,unicode-math}
\setmathfont{Latin Modern Math}
\usepackage{enumitem,needspace,xurl,hyperref}
\setlength{\parindent}{0pt}
\setlength{\parskip}{5pt}
\setlength{\emergencystretch}{3em}
\newcommand{\N}{\mathbb{N}}
\newcommand{\Z}{\mathbb{Z}}
\newcommand{\Q}{\mathbb{Q}}
\newcommand{\R}{\mathbb{R}}
\newcommand{\C}{\mathbb{C}}
\newcommand{\F}{\mathbb{F}}
\newcommand{\A}{\mathbb{A}}
\newcommand{\PP}{\mathbb{P}}
\newcommand{\E}{\mathbb{E}}
\newcommand{\eps}{\varepsilon}
\newcommand{\dd}{\,\mathrm{d}}
\begin{document}
\section*{UnsolvedMath ID: 30006278}
\textbf{Source problem code:} catalogue entry
\subsection{Definitions and mathematical statement}
For a fixed integer \(d\) in the intended range, let \(X_d=\operatorname{Hilb}^{d}(\C^{2})\), the Hilbert scheme parametrizing zero-dimensional subschemes of \(\C^{2}\) of length \(d\). The supplied record uses a torus \(T\), equivariant parameters \(t_1,t_2\), curve classes \(\beta_n\) indexed by \(n\geq 0\), and equivariant cohomology insertions \(a_{\mu}=|\mu\rangle\). It does not specify the torus action, the classes \(\beta_n\), or the indexing and normalization of these insertions. These symbols therefore retain their intended source meanings; the formulation below is conditional on supplying those definitions.

Write \(\overline{\mathcal M}_{g,r}(X_d,\beta_n)\) for the Deligne--Mumford stack of stable maps from connected nodal curves of arithmetic genus \(g\), with \(r\) marked points, to \(X_d\), representing \(\beta_n\). Here the bar and connected-domain convention express the usual stable-map interpretation, which requires confirmation against the source. Let
\[
\operatorname{ev}_i:\overline{\mathcal M}_{g,r}(X_d,\beta_n)\longrightarrow X_d,
\qquad 1\leq i\leq r,
\]
be the evaluation maps. The supplied excerpt states that this stack is equivariantly proper and carries a \(T\)-equivariant virtual fundamental class. Subject to the intended equivariant integration convention, define the primary invariant
\[
I_{g,r,n}(a_{\mu_1},\ldots,a_{\mu_r})
=
\int_{[\overline{\mathcal M}_{g,r}(X_d,\beta_n)]^{\mathrm{vir}}_T}
\prod_{i=1}^{r}\operatorname{ev}_i^{*}(a_{\mu_i})
\]
and its generating series in the formal variable \(q\),
\[
F_{g,r}^{(d)}(\mu_1,\ldots,\mu_r;q)
=
\sum_{n\geq0}q^n I_{g,r,n}(a_{\mu_1},\ldots,a_{\mu_r}).
\]
The source writes this generating series as \(\langle\mu_1,\ldots,\mu_r\rangle_g^{X_d}\) and places it in \(\Q(t_1,t_2,q)\). A formal generating series and a rational function are distinct objects: the rational-function assertion and any expansion convention are retained as source qualifications requiring verification.

\textbf{Problem.} Give a method to compute these primary equivariant Gromov--Witten invariants, or equivalently their generating series, for every integer \(g\geq0\), for the intended values of \(d\), every permitted number \(r\) of markings, and all permitted insertion labels \(\mu_1,\ldots,\mu_r\), including every indexed curve class \(\beta_n\), \(n\geq0\). The ranges of \(d\), \(r\), and the labels, the curve-class definitions, and the treatment of unstable degree-zero cases must be supplied before this becomes a fully specified computational problem. No descendant insertions occur in the displayed source formula.

For context, the excerpt attributes the genus-zero computation to R. Pandharipande and A. Okounkov, citing its reference [5], and says that the genus-zero theory is governed by quantum multiplication
\[
M_D=D\star:
H_T^*(X_d)[[q]]\longrightarrow H_T^*(X_d)[[q]].
\]
Neither \(D\) nor the bibliographic entry [5] is defined in the supplied record.
\subsection{Short English statement}
Find a method computing the displayed primary equivariant Gromov--Witten generating series of \(\operatorname{Hilb}^{d}(\C^2)\) in every genus and for all intended curve classes and insertions. The supplied record leaves essential conventions unspecified, including the torus action, the classes \(\beta_n\), and the insertion basis. Its partially solved status and literature assessment are supplied curation, not independently verified conclusions.
\subsection{Sources}
\begin{itemize}
\item Local UnsolvedMath registry, \texttt{data/unsolvedmath/problems.json}, permanent ID \(30006278\), title ``All-Genus Gromov--Witten Invariants of Hilbert Schemes,'' exact source code \texttt{catalogue entry}. The supplied record identifies the question as Question 3 and gives proposed year 2025; no proposer is named.
\item \emph{Recent Trends in Algebraic Geometry} (2025), Oberwolfach Reports, pp.\ 1451--1516, DOI \texttt{10.4171/owr/2025/28}. Supplied report citation and literature-source label ``original OWR report'': \url{https://doi.org/10.4171/owr/2025/28}.
\item Supplied literature-source label ``You, higher-genus theory of Hilbert schemes'': \url{https://people.math.ethz.ch/~rahul/HilbC2YY.pdf}. A complete bibliographic citation is not supplied.
\item The excerpt attributes the genus-zero computation to R. Pandharipande and A. Okounkov through reference [5]. The corresponding bibliographic entry is absent from the supplied record.
\item Literature assessment dated \texttt{2026-08-22}: the supplied curation labels the problem partially solved, states that the genus-zero theory is known and higher-genus theory has been related to local-curve theories, and reports that no method computing all the stated equivariant invariants in every genus is established. This is supplied dated curation, not an independently verified assessment. The background explicitly requests independent verification of time-sensitive status and cited claims.
\item The registry labels the statement \texttt{corrected\_verified} and reports that the official Oberwolfach report was checked and the notation repaired. These are supplied verification claims; no independent source inspection is asserted here.
\end{itemize}
\end{document}
